Markov 性常被简称为``未来只依赖现在''。真正关键不在遗忘历史，而在状态是否已经包含了预测未来所需的信息。若状态选得太粗，历史仍会通过遗漏变量影响未来，过程便不再 Markov。

这一单元从局部转移逐步走向长期行为：

\begin{enumerate}
\tightlist
\item
  Markov 性质、转移矩阵与多步演化：一步规则怎样通过矩阵乘法生成整个路径分布。
\item
  可达、通信类与周期性：先用有向图看链被分成哪些区域，为什么有些链会永久振荡。
\item
  常返、暂留与返回时间：状态会不会最终回来，以及有限链与无限链为何不同。
\item
  平稳分布、可逆性与遍历收敛：区分``不随时间改变的分布''\,``从任意初值趋向的分布''和``单条路径的长期频率''。
\item
  首步分析、击中时间与吸收链：把未来问题按第一步拆开，计算到达概率、等待时间和访问次数。
\end{enumerate}

这五节的主线是：局部转移规则并不足以直接说明长期稳定；通信结构、常返性和周期性必须先被辨认。

举一个具体例子。设

\[
P=
\begin{bmatrix}
0.8&0.2\\
0.4&0.6
\end{bmatrix}
\]
是简化的每日天气链。晴天的转移行为是 $0.8$ 概率保持晴、$0.2$ 转雨；雨天则是 $0.4$ 转晴、$0.6$ 保持雨。解 $\pi P=\pi$ 得平稳分布 $\pi=(2/3,\,1/3)$。它不可约（两状态互通）、有自环故非周期，因此无论从晴天还是雨天出发，$P^n$ 都收敛到 $\pi$。一个关键测试：从雨天出发，两天后回雨的概率是多少？$P^2$ 的右下角是 $(0.4)(0.2)+(0.6)(0.6)=0.44$——比一步的 $0.6$ 下降了，正在往平稳值 $1/3\approx0.33$ 靠拢。这个链虽小，却包含了本单元所有核心概念。

\subsection{怎么读这一章}

先记住：Markov 性取决于状态表示；平稳分布是转移算子的固定点；逐时刻分布收敛与单条路径时间平均是不同结论。遇到``链最终会稳定''之类说法时，应追问不可约、正常返与非周期条件是否齐全。

若要实际计算，再学习转移矩阵幂、通信类、首步分析和吸收链基本矩阵；若要进入 MCMC、PageRank 或强化学习，则应进一步掌握可逆性、谱隙与遍历定理。返回性分类的证明可在需要处理无限状态链时再深入。

\subsection{本章篇目}

以下篇目按概念依赖排列；若从具体问题进入，也可以先打开最接近当前困惑的一篇，再沿篇内链接回补。

\begin{itemize}
\tightlist
\item \hyperref[sec:06-Random-Processes/03-Markov-Chains/01]{``Markov 性质、转移矩阵与多步演化''}；
\item \hyperref[sec:06-Random-Processes/03-Markov-Chains/02]{``可达、通信类与周期性''}；
\item \hyperref[sec:06-Random-Processes/03-Markov-Chains/03]{``常返、暂留与返回时间''}；
\item \hyperref[sec:06-Random-Processes/03-Markov-Chains/04]{``平稳分布、可逆性与遍历收敛''}；
\item \hyperref[sec:06-Random-Processes/03-Markov-Chains/05]{``首步分析、击中时间与吸收链''}。
\end{itemize}
